\documentclass[
    language=english,
    doctype=recipe-book,
    institution=none,
]{../../omnilatex}

\RequirePackage{config/document-settings}

\begin{document}
\maketitle

\frontmatter

\chapter*{Introduction}

This book is a love letter to the food I grew up eating. My grandmother cooked the way she did everything else---with precision, generosity, and an unwillingness to accept anything less than the best. Her kitchen was the heart of the house, and the food she made there was the heart of the family.

I've spent the last ten years trying to recreate her recipes, and what I've learned is this: the recipes are the easy part. The hard part is the attention. My grandmother didn't measure ingredients by weight or volume. She measured them by feel, by sight, by the sound of sizzling butter in a pan. She didn't follow timers. She followed her instincts, developed over sixty years of daily cooking.

I can't teach you her instincts. But I can give you her recipes, and I can tell you what I've learned about the principles behind them---the why, not just the what. If you understand why a recipe works, you can adapt it, improve it, and eventually make it your own.

These recipes are not all from my grandmother. Some are from my travels, some from friends, some from cookbooks that changed my life. But they all share a philosophy: cook with good ingredients, cook with care, and cook for people you love.

--- Maria Santos, June 2024

\tableofcontents

\mainmatter

\chapter{Breakfast}

\section*{Classic French Toast with Caramelized Bananas}

\begin{recipebox}
\textbf{Prep Time:} 10 minutes\\
\textbf{Cook Time:} 15 minutes\\
\textbf{Servings:} 4\\
\textbf{Difficulty:} Easy
\end{recipebox}

\subsection*{Ingredients}

\begin{multicols}{2}
\begin{itemize}[leftmargin=*]
    \item 8 slices thick-cut brioche (day-old is best)
    \item 4 large eggs
    \item 1 cup whole milk
    \item 1 teaspoon vanilla extract
    \item 1 teaspoon ground cinnamon
    \item 1/4 teaspoon ground nutmeg
    \item Pinch of salt
    \item 2 tablespoons unsalted butter
    \item 2 ripe bananas, sliced
    \item 2 tablespoons brown sugar
    \item 1 tablespoon dark rum (optional)
    \item Powdered sugar for dusting
    \item Maple syrup for serving
\end{itemize}
\end{multicols}

\subsection*{Instructions}

\begin{enumerate}
    \item In a shallow dish, whisk together eggs, milk, vanilla, cinnamon, nutmeg, and salt until smooth.
    \item Heat a large skillet or griddle over medium heat. Melt 1 tablespoon of butter.
    \item Dip each slice of brioche in the egg mixture, allowing it to soak for 15--20 seconds per side. Don't over-soak or the bread will fall apart.
    \item Cook the French toast for 3--4 minutes per side, until golden brown and crispy at the edges. Add more butter as needed between batches.
    \item While the last batch cooks, melt remaining butter in a separate skillet. Add banana slices and brown sugar. Cook for 2--3 minutes, stirring gently, until bananas are caramelized.
    \item If using rum, carefully add it to the pan and ignite (or let it cook off for 30 seconds).
    \item Serve French toast topped with caramelized bananas, a dusting of powdered sugar, and maple syrup.
\end{enumerate}

\subsection*{Cook's Notes}
\begin{quote}
The secret to good French toast is the bread. Brioche is ideal because it's rich and eggy, but challah works too. The key is to use bread that's at least a day old---fresh bread gets soggy. If your bread is fresh, toast it lightly before dipping.

Also: don't press down on the French toast while it's cooking. Let it cook undisturbed. The crust forms in the first two minutes, and if you mess with it, you'll破坏 the crust.
\end{quote}

\section*{Granola with Honey and Toasted Almonds}

\begin{recipebox}
\textbf{Prep Time:} 10 minutes\\
\textbf{Cook Time:} 25 minutes\\
\textbf{Servings:} 8\\
\textbf{Difficulty:} Easy
\end{recipebox}

\subsection*{Ingredients}

\begin{itemize}
    \item 3 cups old-fashioned rolled oats
    \item 1 cup raw almonds, roughly chopped
    \item 1/2 cup raw pecans, roughly chopped
    \item 1/2 cup unsweetened coconut flakes
    \item 1/4 cup sunflower seeds
    \item 1/4 cup pumpkin seeds
    \item 1/3 cup honey
    \item 1/4 cup coconut oil, melted
    \item 1 teaspoon vanilla extract
    \item 1 teaspoon ground cinnamon
    \item 1/2 teaspoon salt
    \item 1 cup dried cranberries (added after baking)
\end{itemize}

\subsection*{Instructions}

\begin{enumerate}
    \item Preheat oven to 325°F (165°C). Line a large baking sheet with parchment paper.
    \item In a large bowl, combine oats, almonds, pecans, coconut flakes, sunflower seeds, and pumpkin seeds.
    \item In a small bowl, whisk together honey, melted coconut oil, vanilla, cinnamon, and salt.
    \item Pour the wet ingredients over the dry ingredients and stir until everything is evenly coated.
    \item Spread the mixture in an even layer on the prepared baking sheet.
    \item Bake for 20--25 minutes, stirring every 10 minutes, until golden brown and fragrant.
    \item Remove from oven and let cool completely on the baking sheet (it will crisp up as it cools).
    \item Once cool, stir in dried cranberries. Store in an airtight container for up to 2 weeks.
\end{enumerate}

\subsection*{Cook's Notes}
\begin{quote}
Granola is infinitely adaptable. Don't like cranberries? Use raisins or dried cherries. Don't like almonds? Use walnuts or hazelnuts. The only thing you really need to get right is the ratio of wet to dry ingredients. If your granola is too dry, it won't clump. If it's too wet, it won't crisp. Start with the recipe as written and adjust next time based on your results.
\end{quote}

\chapter{Lunch}

\section*{Mediterranean Grain Bowl}

\begin{recipebox}
\textbf{Prep Time:} 15 minutes\\
\textbf{Cook Time:} 20 minutes\\
\textbf{Servings:} 4\\
\textbf{Difficulty:} Easy
\end{recipebox}

\subsection*{Ingredients}

\begin{multicols}{2}
\begin{itemize}[leftmargin=*]
    \item 1 1/2 cups quinoa, rinsed
    \item 1 can (15 oz) chickpeas, drained and rinsed
    \item 2 cups cherry tomatoes, halved
    \item 1 English cucumber, diced
    \item 1/2 cup Kalamata olives, halved
    \item 1/4 cup red onion, thinly sliced
    \item 1/2 cup crumbled feta cheese
    \item 1/4 cup fresh parsley, chopped
    \item 2 tablespoons fresh mint, chopped
    \item 1/4 cup extra virgin olive oil
    \item 2 tablespoons red wine vinegar
    \item 1 tablespoon lemon juice
    \item 1 clove garlic, minced
    \item 1 teaspoon dried oregano
    \item Salt and pepper to taste
\end{itemize}
\end{multicols}

\subsection*{Instructions}

\begin{enumerate}
    \item Cook quinoa according to package directions. Let cool slightly.
    \item In a small bowl, whisk together olive oil, red wine vinegar, lemon juice, garlic, and oregano.
    \item In a large bowl, combine cooled quinoa, chickpeas, tomatoes, cucumber, olives, and red onion.
    \item Pour dressing over the salad and toss to combine.
    \item Season with salt and pepper.
    \item Divide among bowls and top with feta, parsley, and mint.
\end{enumerate}

\subsection*{Cook's Notes}
\begin{quote}
This bowl is infinitely better if you marinate the chickpeas. Drain and rinse them, toss with a tablespoon of olive oil, a pinch of cumin, a pinch of paprika, and a squeeze of lemon. Roast at 400°F for 20 minutes until crispy. The crunch changes everything.

Also, don't skip the mint. It seems like a small thing, but it makes a huge difference.
\end{quote}

\section*{Classic Tomato Soup}

\begin{recipebox}
\textbf{Prep Time:} 10 minutes\\
\textbf{Cook Time:} 40 minutes\\
\textbf{Servings:} 6\\
\textbf{Difficulty:} Easy
\end{recipebox}

\subsection*{Ingredients}

\begin{itemize}
    \item 2 tablespoons olive oil
    \item 1 large onion, diced
    \item 3 cloves garlic, minced
    \item 2 cans (28 oz each) San Marzano tomatoes
    \item 2 cups vegetable broth
    \item 1 teaspoon sugar
    \item 1 teaspoon dried basil
    \item 1/2 teaspoon dried oregano
    \item Salt and pepper to taste
    \item 1/2 cup heavy cream (optional)
    \item Fresh basil for garnish
\end{itemize}

\subsection*{Instructions}

\begin{enumerate}
    \item Heat olive oil in a large pot over medium heat. Add onion and cook until soft and translucent, about 5 minutes.
    \item Add garlic and cook for 1 minute until fragrant.
    \item Add tomatoes (with their juices), broth, sugar, basil, and oregano. Bring to a boil, then reduce heat and simmer for 30 minutes.
    \item Remove from heat and let cool slightly. Blend with an immersion blender (or in batches in a regular blender) until smooth.
    \item Return to heat if needed. Stir in cream if using. Season with salt and pepper.
    \item Serve hot, garnished with fresh basil and a drizzle of olive oil.
\end{enumerate}

\subsection*{Cook's Notes}
\begin{quote}
The secret to good tomato soup is San Marzano tomatoes. They're sweeter and less acidic than regular canned tomatoes. If you can't find them, add an extra teaspoon of sugar.

This soup is also excellent with a grilled cheese sandwich, but that's a different recipe.
\end{quote}

\chapter{Dinner}

\section*{Slow-Braised Short Ribs}

\begin{recipebox}
\textbf{Prep Time:} 20 minutes\\
\textbf{Cook Time:} 3 hours\\
\textbf{Servings:} 4\\
\textbf{Difficulty:} Moderate
\end{recipebox}

\subsection*{Ingredients}

\begin{multicols}{2}
\begin{itemize}[leftmargin=*]
    \item 4 lbs bone-in beef short ribs
    \item Salt and freshly ground black pepper
    \item 2 tablespoons vegetable oil
    \item 2 large onions, diced
    \item 3 carrots, peeled and cut into 1-inch pieces
    \item 4 cloves garlic, minced
    \item 2 tablespoons tomato paste
    \item 1/4 cup all-purpose flour
    \item 2 cups dry red wine
    \item 2 cups beef broth
    \item 2 sprigs fresh thyme
    \item 2 sprigs fresh rosemary
    \item 2 bay leaves
    \item 1 tablespoon balsamic vinegar
\end{itemize}
\end{multicols}

\subsection*{Instructions}

\begin{enumerate}
    \item Preheat oven to 325°F (165°C).
    \item Season short ribs generously with salt and pepper.
    \item Heat oil in a Dutch oven over medium-high heat. Working in batches, brown short ribs on all sides, about 3--4 minutes per side. Transfer to a plate.
    \item Add onions and carrots to the pot. Cook until onions are soft and golden, about 8 minutes.
    \item Add garlic and tomato paste. Cook for 2 minutes, stirring constantly.
    \item Sprinkle flour over the vegetables and stir for 1 minute.
    \item Pour in red wine, scraping up any browned bits from the bottom of the pot. Let simmer for 3 minutes.
    \item Add broth, thyme, rosemary, and bay leaves. Return short ribs to the pot.
    \item Bring to a simmer, cover, and transfer to oven. Braise for 2.5--3 hours, until meat is fork-tender.
    \item Remove short ribs. Discard herb sprigs and bay leaves. Stir in balsamic vinegar.
    \item If sauce is too thin, simmer on stovetop until thickened.
    \item Serve short ribs over mashed potatoes or polenta, spooning sauce over the top.
\end{enumerate}

\subsection*{Cook's Notes}
\begin{quote}
The key to braising is patience. Don't rush it. The meat should be so tender that it falls off the bone with a gentle push of a fork. If it's not that tender, it's not done.

Also, the wine matters. Don't cook with wine you wouldn't drink. A good Côtes du Rhône is perfect here.
\end{quote}

\section*{Lemon Herb Roasted Chicken}

\begin{recipebox}
\textbf{Prep Time:} 20 minutes\\
\textbf{Cook Time:} 1 hour 15 minutes\\
\textbf{Servings:} 4--6\\
\textbf{Difficulty:} Easy
\end{recipebox}

\subsection*{Ingredients}

\begin{itemize}
    \item 1 whole chicken (4--5 lbs)
    \item 2 lemons
    \item 1 head of garlic
    \item Fresh herbs (thyme, rosemary, sage)
    \item 3 tablespoons olive oil
    \item 1 teaspoon salt
    \item 1/2 teaspoon pepper
    \item 1 pound baby potatoes, halved
    \item 1 pound carrots, peeled and cut into chunks
    \item 1 large onion, quartered
\end{itemize}

\subsection*{Instructions}

\begin{enumerate}
    \item Preheat oven to 425°F (220°C).
    \item Pat chicken dry with paper towels. Rub with 2 tablespoons olive oil, salt, and pepper.
    \item Cut one lemon in half and stuff into the chicken cavity along with half the garlic and a bundle of herbs.
    \item Place chicken in a large roasting pan. Arrange potatoes, carrots, and onion around the chicken. Drizzle remaining olive oil over vegetables.
    \item Cut remaining lemon into wedges and scatter around the pan. Squeeze remaining garlic cloves over everything.
    \item Roast for 1 hour 15 minutes, or until internal temperature reaches 165°F (74°C) and juices run clear.
    \item Let rest for 15 minutes before carving. Serve with vegetables and pan juices.
\end{enumerate}

\subsection*{Cook's Notes}
\begin{quote}
The secret to a crispy-skinned roasted chicken is dry skin. Pat it dry, really dry, before you oil it. Moisture is the enemy of crispiness.

Also, don't baste. I know every recipe tells you to baste. Don't do it. Every time you open the oven, you lose heat, and the skin gets soggy. Trust the oven. Let it do its job.
\end{quote}

\chapter{Dessert}

\section*{Classic Tiramisu}

\begin{recipebox}
\textbf{Prep Time:} 30 minutes\\
\textbf{Cook Time:} None (chilling time: 4 hours)\\
\textbf{Servings:} 8--10\\
\textbf{Difficulty:} Moderate
\end{recipebox}

\subsection*{Ingredients}

\begin{multicols}{2}
\begin{itemize}[leftmargin=*]
    \item 6 egg yolks
    \item 3/4 cup sugar
    \item 1 1/4 cups mascarpone cheese
    \item 2 cups heavy cream
    \item 2 cups strong espresso, cooled
    \item 3 tablespoons coffee liqueur (optional)
    \item 2 packages ladyfinger cookies (Savoiardi)
    \item Unsweetened cocoa powder for dusting
    \item Dark chocolate shavings for garnish
\end{itemize}
\end{multicols}

\subsection*{Instructions}

\begin{enumerate}
    \item In a large bowl, beat egg yolks and sugar until thick and pale yellow, about 4 minutes.
    \item Add mascarpone and beat until smooth.
    \item In a separate bowl, whip heavy cream to stiff peaks. Fold into mascarpone mixture.
    \item Combine espresso and coffee liqueur (if using) in a shallow dish.
    \item Briefly dip each ladyfinger in the espresso (don't soak---just a quick dip). Arrange in a single layer in a 9x13 inch dish.
    \item Spread half the mascarpone mixture over the ladyfingers.
    \item Repeat with another layer of dipped ladyfingers and remaining mascarpone.
    \item Cover and refrigerate for at least 4 hours, preferably overnight.
    \item Dust generously with cocoa powder and top with chocolate shavings before serving.
\end{enumerate}

\subsection*{Cook's Notes}
\begin{quote}
The most common mistake with tiramisu is soaking the ladyfingers too long. They should be wet on the outside but still firm in the center. If they get soggy, the whole thing collapses.

Also, it tastes better the next day. Make it the day before you plan to serve it.
\end{quote}

\section*{Chocolate Lava Cake}

\begin{recipebox}
\textbf{Prep Time:} 15 minutes\\
\textbf{Cook Time:} 12 minutes\\
\textbf{Servings:} 4\\
\textbf{Difficulty:} Moderate
\end{recipebox}

\subsection*{Ingredients}

\begin{itemize}
    \item 4 oz bittersweet chocolate (70\% cacao)
    \item 1/2 cup (1 stick) unsalted butter
    \item 2 large eggs
    \item 2 large egg yolks
    \item 1/4 cup sugar
    \item Pinch of salt
    \item 2 tablespoons all-purpose flour
    \item Butter and cocoa powder for ramekins
    \item Vanilla ice cream for serving
\end{itemize}

\subsection*{Instructions}

\begin{enumerate}
    \item Preheat oven to 425°F (220°C). Butter four 6-ounce ramekins and dust with cocoa powder.
    \item Melt chocolate and butter together in a double boiler or microwave. Let cool slightly.
    \item In a large bowl, whisk eggs, egg yolks, sugar, and salt until thick and pale, about 3 minutes.
    \item Fold chocolate mixture into egg mixture. Gently fold in flour.
    \item Divide batter among prepared ramekins.
    \item Bake for 12--14 minutes until edges are firm but center is soft.
    \item Let cool for 1 minute, then invert onto plates.
    \item Serve immediately with vanilla ice cream.
\end{enumerate}

\subsection*{Cook's Notes}
\begin{quote}
Timing is everything. The difference between a perfect lava cake and a regular chocolate cake is about 90 seconds. Set a timer. Check at 12 minutes. The edges should be set, but the center should jiggle slightly when you tap the ramekin.

If you're nervous, make a practice one first. Or make extra batter and bake them in batches. The first one is always a learning experience.
\end{quote}

\section*{Nutrition Information}

\begin{table}[h]
\centering
\small
\begin{tabular}{lrrr}
\toprule
\textbf{Recipe} & \textbf{Calories} & \textbf{Protein} & \textbf{Prep Time} \\
\midrule
French Toast & 420 & 12g & 10 min \\
Granola & 280 & 8g & 10 min \\
Grain Bowl & 380 & 15g & 15 min \\
Tomato Soup & 220 & 6g & 10 min \\
Short Ribs & 580 & 35g & 20 min \\
Roasted Chicken & 450 & 38g & 20 min \\
Tiramisu & 480 & 8g & 30 min \\
Lava Cake & 520 & 9g & 15 min \\
\bottomrule
\end{tabular}
\caption{Approximate nutritional information per serving}
\end{table}

\end{document}
